% Section 10 Header
\section{Section 10 — Fast \& Slow Pointers}

\begin{frame}[c]{\empty}
    \begin{center}
        \Large\textbf{\secname}
    \end{center}
\end{frame}

% Slide 1: Concept
\begin{frame}{What is the Fast \& Slow Pointers Pattern?}
  \begin{itemize}
    \item \textbf{The Core Idea:} Use two pointers traversing a linear data structure (linked list/array) at different speeds.
    \item \textbf{Why it works:} If there is a cycle, the fast pointer (moving 2 steps at a time) will eventually meet the slow pointer (moving 1 step at a time) in $O(N)$ time and $O(1)$ space.
    \item \textbf{Common Triggers:}
      \begin{itemize}
        \item Cycle detection in a linked list.
        \item Finding the starting node of a cycle.
        \item Finding the middle node of a list (when the fast pointer reaches the end, the slow pointer is in the middle).
        \item Identifying duplicates in structures where values can be treated as indices.
      \end{itemize}
  \end{itemize}
\end{frame}

\begin{frame}{Cycle Detection Mechanics}
  \begin{itemize}
    \item \textbf{Initialize:} `slow = head`, `fast = head`.
    \item \textbf{Advance:}
      \begin{itemize}
        \item `slow = slow.next`
        \item `fast = fast.next.next`
      \end{itemize}
    \item \textbf{Check:} If `slow == fast`, a cycle is detected! If `fast` or `fast.next` is null, there is no cycle.
    \item \textbf{Cycle Start (Floyd's algorithm):} Once met, reset `slow` to `head`. Move both pointers 1 step at a time. The point where they meet again is the start of the cycle.
  \end{itemize}
\end{frame}

% Problem 1: Linked List Cycle
\begin{frame}[c]{\empty}
    \begin{center}
        \Large \textbf{Problem 1: \href{https://leetcode.com/problems/linked-list-cycle/}{Linked List Cycle}}
    \end{center}
\end{frame}

% Problem 2: Middle of the Linked List
\begin{frame}[c]{\empty}
    \begin{center}
        \Large \textbf{Problem 2: \href{https://leetcode.com/problems/middle-of-the-linked-list/}{Middle of the Linked List}}
    \end{center}
\end{frame}

% Problem 3: Find the Duplicate Number
\begin{frame}[c]{\empty}
    \begin{center}
        \Large \textbf{Problem 3: \href{https://leetcode.com/problems/find-the-duplicate-number/}{Find the Duplicate Number}}
    \end{center}
\end{frame}
